%-------------------------
% Resume in Latex
% Author : Jake Gutierrez
% Based off of: https://github.com/sb2nov/resume
% License : MIT
%------------------------

\documentclass[letterpaper,11pt]{article}

\usepackage{latexsym}
\usepackage[empty]{fullpage}
\usepackage{titlesec}
\usepackage{marvosym}
\usepackage[usenames,dvipsnames]{color}
\usepackage{verbatim}
\usepackage{enumitem}
\usepackage[hidelinks]{hyperref}
\usepackage{fancyhdr}
\usepackage[english]{babel}
\usepackage{tabularx}
\usepackage{fontawesome5}
\usepackage{multicol}
\setlength{\multicolsep}{-3.0pt}
\setlength{\columnsep}{-1pt}
\input{glyphtounicode}

%----------FONT OPTIONS----------
% sans-serif
% \usepackage[sfdefault]{FiraSans}
% \usepackage[sfdefault]{roboto}
% \usepackage[sfdefault]{noto-sans}
% \usepackage[default]{sourcesanspro}

% serif
% \usepackage{CormorantGaramond}
% \usepackage{charter}

\pagestyle{fancy}
\fancyhf{} % clear all header and footer fields
\fancyfoot{}
\renewcommand{\headrulewidth}{0pt}
\renewcommand{\footrulewidth}{0pt}

% Adjust margins
\addtolength{\oddsidemargin}{-0.6in}
\addtolength{\evensidemargin}{-0.5in}
\addtolength{\textwidth}{1.19in}
\addtolength{\topmargin}{-.7in}
\addtolength{\textheight}{1.4in}

\urlstyle{same}

\raggedbottom
\raggedright
\setlength{\tabcolsep}{0in}

% Sections formatting
\titleformat{\section}{
  \vspace{-4pt}\scshape\raggedright\large\bfseries
}{}{0em}{}[\color{black}\titlerule \vspace{-5pt}]

% Ensure that generate pdf is machine readable/ATS parsable
\pdfgentounicode=1

%-------------------------
% Custom commands
\newcommand{\resumeItem}[1]{
\item\small{
    {#1 \vspace{-2pt}}
  }
}

\newcommand{\classesList}[4]{
\item\small{
    {#1 #2 #3 #4 \vspace{-2pt}}
  }
}

\newcommand{\resumeSubheading}[4]{
  \vspace{-2pt}
\item
  \begin{tabular*}{1.0\textwidth}[t]{l@{\extracolsep{\fill}}r}
    \textbf{#1} & \textbf{\small #2} \\
    \textit{\small#3} & \textit{\small #4} \\
  \end{tabular*}\vspace{-7pt}
}

\newcommand{\resumeSubSubheading}[2]{
\item
  \begin{tabular*}{0.97\textwidth}{l@{\extracolsep{\fill}}r}
    \textit{\small#1} & \textit{\small #2} \\
  \end{tabular*}\vspace{-7pt}
}

\newcommand{\resumeProjectHeading}[2]{
\item
  \begin{tabular*}{1.001\textwidth}{l@{\extracolsep{\fill}}r}
    \small#1 & \textbf{\small #2}\\
  \end{tabular*}\vspace{-7pt}
}

\newcommand{\resumeSubItem}[1]{\resumeItem{#1}\vspace{-4pt}}

\renewcommand\labelitemi{$\vcenter{\hbox{\tiny$\bullet$}}$}
\renewcommand\labelitemii{$\vcenter{\hbox{\tiny$\bullet$}}$}

\newcommand{\resumeSubHeadingListStart}{
\begin{itemize}[leftmargin=0.0in, label={}]}
    \newcommand{\resumeSubHeadingListEnd}{
\end{itemize}}
\newcommand{\resumeItemListStart}{
\begin{itemize}}
    \newcommand{\resumeItemListEnd}{
\end{itemize}\vspace{-5pt}}

%-------------------------------------------
%%%%%%  RESUME STARTS HERE  %%%%%%%%%%%%%%%%%%%%%%%%%%%%

\begin{document}

\begin{center}
  {\Huge \scshape Ankit Khushwaha} \\ \vspace{2pt}
  \raisebox{-0.2\height}\faPhone\ \href{tel:+917428061540}{\underline{+91-7428061540}} ~
  \href{mailto:ankitkhushwaha.dev@gmail.com}{\raisebox{-0.2\height}\faEnvelope\
  \underline{ankitkhushwaha.dev@gmail.com}} ~
  \href{https://www.linkedin.com/in/ankit-khushwaha-2073a0252/}{\raisebox{-0.2\height}\faLinkedin\ \underline{linkedin}}
  ~
  \href{https://github.com/ankitkhushwaha}{\raisebox{-0.2\height}\faGithub\ \underline{github}}
  ~
  \href{https://ankitkdev.com/}{\raisebox{-0.2\height}\faGlobe\ \underline{blog}}

  \vspace{-8pt}
\end{center}

%-----------EDUCATION-----------
\section{Education}
\resumeSubHeadingListStart
\resumeSubheading
{Indian Institute Of Technology Dharwad}{Dec. 2022 -- Present}
{Bachelor of Science-Master of Science, Physics}{Dharwad, Karnataka}
\resumeSubHeadingListEnd
\vspace{-16pt}

%-----------PROGRAMMING SKILLS-----------

\section{Technical Skills}
\begin{itemize}[leftmargin=0.15in, label={}]
    \small{
    \item \textbf{Languages:} C, C++, Python\\
      \vspace{-2pt}
    \item \textbf{Systems \& Tools:} Linux, QEMU, U-boot, GDB, GNU Make, Kbuild, Git \\
      \vspace{-2pt}
    \item \textbf{Core Skills:} Linux Device Driver, Kernel Modules,
 Kernel Debugging, Device Tree, System Programming \\

      \vspace{-2pt}
    \item \textbf{Kernel Subsystems:} I2C, Platform/Custom Bus, Character Drivers, Sysfs, Procfs, Timers, Workqueues, Interrupt Handling, Synchronization \\

      \vspace{-2pt}
    \item \textbf{Familiar With:}  Shell Scripting, Computer Architecture, Web Development\\
    }
\vspace{-2pt}
\end{itemize}
\vspace{-18pt}

%--------------------------------
\section{Experience}
\resumeSubHeadingListStart
\vspace{-2pt}
\resumeSubheading
{Linux Foundation Mentorship (LFX '25) - Linux Kernel}{Sept 2025 - Dec 2025}
{Mentee}{\href{https://git.kernel.org/pub/scm/linux/kernel/git/torvalds/linux.git/log/?qt=grep\&q=ankitkhushwaha.linux\%40gmail.com}{Upstream
Linux Kernel Patches}}
\vspace{-6pt}
\resumeItemListStart
\resumeItem{Contributed 10+ patches across Linux kernel selftests in
net, mm, tls, mptcp, and user\_events subsystems.}
\resumeItem{Analyzed syzbot bugs using QEMU and
analyzed execution paths with ftrace \& gdb to identify root causes.}
\resumeItem{Performed memory error analysis using KASAN to detect and
fix uninitialized variables and invalid memory usage.}
\resumeItem{Collaborated with mentors and upstream contributors through mailing-list discussions, code reviews, \& patch revisions.}

\resumeItemListEnd
\resumeSubHeadingListEnd
\vspace{-25pt}

\resumeSubHeadingListStart

\resumeSubheading
{Acadia Ventures}{July 2025 - Aug 2025}
{Research and Development Intern}{Remote}
\resumeItemListStart
\resumeItem{Built a structured testing framework with fixtures to
  manage inter-test dependencies,
streamlining test maintenance and enabling faster debugging cycles.}
\resumeItem{Implemented secure JWT-based authentication with
role-based access, strengthening API security and control.}
\resumeItem{Built notification system using event
triggers and async tasks, improving user engagement with timely alerts.}
\resumeItemListEnd

\resumeSubHeadingListEnd
\vspace{-16pt}

% -----------EXPERIENCE-----------

% \section{Open-Source Contribution}
% \resumeSubHeadingListStart
% \resumeSubheading
% {Scientific Python Libraries}
% {{\href{https://github.com/sunpy/sunpy/issues?q=is\%3Apr+is\%3Aopen+author\%3Aankitkhushwaha}{SunPy}}
%  -
% \href{https://github.com/StingraySoftware/stingray/issues?q=is\%3Apr+is\%3Aopen+author\%3Aankitkhushwaha}{Stingray}}
% {}{}
% \vspace{-20pt}
% \resumeItemListStart
% \resumeItem{Contributed to \textbf{SunPy} and \textbf{Stingray}, open-source Python libraries for scientific data analysis.}
% \resumeItem{Enhanced time-series analysis, data processing, ingestion, and visualization capabilities across both libraries.}
% \resumeItem{Redesigned documentation UI and Improved data handling, utility functions, documentation, and maintainability.}
% \resumeItem{Contributed to CI/CD automation, release workflows, code reviews, and issue discussions with project maintainers.}
% \resumeItemListEnd
% \resumeSubHeadingListEnd
% \vspace{-12pt}

%-----------PROJECTS-----------
\section{Projects}
\vspace{-3pt}

\resumeSubHeadingListStart
\resumeSubheading
{Linux I2C EEPROM Driver}
{\href{https://github.com/ankitkhushwaha/linux-i2c-eeprom-driver}{\faGithub}}
{C, Linux Kernel, I2C, Device Driver}{}
\vspace{-5pt}
\resumeItemListStart
\resumeItem{Developed a Linux \textbf{I2C client driver} for an EEPROM device, implementing \textbf{probe/remove} lifecycle and device--driver matching via the \textbf{I2C ID table}.}
\resumeItem{Implemented byte-level and page-aware \textbf{read/write} access over the I2C bus using the kernel \textbf{I2C transfer APIs}.}
\resumeItem{Exposed EEPROM contents to user space through a \textbf{sysfs/character device interface} with proper \textbf{bounds checking and error handling}.}
\resumeItem{Protected concurrent accesses with \textbf{mutex-based locking} and handled write-cycle delays and transfer failures for reliable operation.}
\resumeItem{Validated the driver with \textbf{beaglebone black} and debugged using kernel logs and \textbf{dynamic debug}.}
\resumeItemListEnd
\vspace{-8pt}
\resumeSubHeadingListEnd

\resumeSubHeadingListStart
\vspace{-3pt}
\resumeSubheading
{Mini SoC Device Model}
{\href{https://github.com/ankitkhushwaha/mini-soc-device-model}{\faGithub}}
{C, Linux Kernel, Device Model, Makefile}{}
\vspace{-3pt}
\resumeItemListStart
\resumeItem{Built a simulated \textbf{SoC sensor hub} as three kernel modules (bus, device, driver) implementing a custom \textbf{\texttt{bus\_type}} modelled on the Linux \textbf{platform bus}.}
\resumeItem{Implemented \textbf{device/driver matching, probe/remove, and sysfs unbinding}, along with \textbf{\texttt{struct device}} lifecycle management .}
\resumeItem{Added \textbf{Device Tree} support by binding \textbf{\texttt{of\_node}} to devices with \texttt{device\_set\_node()} and populating and tearing down devices from DT nodes}
\resumeItem{Tested on \textbf{ARM QEMU} with a cross-compiled kernel and debugged with \textbf{gdb} using kernel GDB scripts.}\resumeItemListEnd
\vspace{-8pt}
\resumeSubHeadingListEnd


\section{Certifications}
\resumeSubHeadingListStart
\resumeSubheading
{The Linux Foundation}
{Issued Aug 2025}
{LFD103: A Beginner's Guide to Linux Kernel Development}
{\href{https://www.credly.com/badges/84d27b87-bf04-43fb-971b-6484b78cb651/public_url}{Credly}}
\vspace{-4pt}
\resumeSubheading
{Udemy}
{Issued July 2026}
{LDD1: Linux Device Driver Programming With Beaglebone Black}
{\href{https://www.udemy.com/certificate/UC-5013297e-6908-4e59-8aef-18bf2e4a293e/}{Udemy}}

\resumeSubHeadingListEnd
\end{document}
